\chapter{Conclusion}
\label{sec:conclusion}
\section{Summary}
In this work, we propose and develop a process modeling tool, especially for the needs process modeling from documents. We explore two possibilties 

Here we summarize the research questions:

RQ1 How do existing tools supported AI-assisted process modeling with documents?

As elaborated in Chapter02, the current tools in research / academic area, mainly focus on the quality of the generated model, accoding to the provided description. To our best knowledage, there is no tools right now to support model auto-generation from document, therefore, to do so an export has to copy and paste text themselves

RQ2 How can an integrated system be designed to support process modeling from documents?
RQ2 Which layout and properties are beneficial for an integrated system that supports AI assisted process modeling with multiple documents as input?
- ?

RQ3 How do user perceive the benefits and limitation of the system?
RQ3 What are advantages and disadvantages of such a system? 

Perceived benefits: Intuitive and reduce the tool switching time..., versioning
Perceived limitations: More features is expected

Participants found the system particularly valuable for creating and exploring multiple models from documents using features like text-model link and process-subprocess model link. However, for simpler tasks, such as building a model from a single document, the system offers limited advantages compared to existing tools.

In summary, we extend the functionality of LLM-based modeling processing to in two area:
1. we offered a new input modality for LLM-based modeling - the input text is selected from the document.
2. a traceability.
3. 


\section{Future Work}
Future work includes:
- Improve and Extend the current features. e.g. the rendering issue , improve ..., text search,,, or support more documents type by leveraging other types... For such a system a lot of stuff can be done.
- Integrated some document intelligent from simple e.g. using AI-to summarize the content, to advanced like chatbot query
- Fine-tuning the LLM for more precise text extraction, as discussed, since one of the key features is the traceablity, a dedicated for such case can be offered.


I. Research Implications
 

 
II. Future Research Directions
 
Future research can be conducted in depth across four areas: user interaction, system functionality, intelligent assistance, and scenario implementation.
First, optimize user-intuitive interaction and visualization design, simplify operational steps, and improve multi-format document rendering capabilities. Clearly distinguish model elements and version relationships through visual cues to enhance overall usability.
Second, establish a refined model management and traceability system. Develop features such as version comparison, fine-grained text-to-model traceability, and intelligent document information filtering to enhance full-lifecycle modeling support capabilities.
Third, research explainable AI-assisted modeling mechanisms. Clarify AI input rules and reasoning processes to enable predictable modification outcomes and controllable workflows, ensuring user autonomy while advancing intelligence levels.
Fourth, integrate search-enhanced generation technology with conversational interaction to build cross-document contextual modeling capabilities. Develop automated model quality assessment and optimization functions to further improve modeling efficiency and quality.
Fifth, conduct adaptation research for practical scenarios such as complex business processes and multi-document collaboration. Promote the integration of the tool with existing systems to establish standardized application paradigms, thereby enhancing the tool’s engineering implementation capabilities and expanding its industry applicability.